% GENERATED FILE. Edit canonical lessons, then run npm run generate:books.
% canonical-source-hash: fnv1a64:9c142d21b1b505fb
% canonical-lessons: BN-C462-olos, BN-C462-kormoth, BN-C462-elomelo, BN-C462-phaki-deoya, BN-C462-gorimosi-kora

\chapter{Lazy and Industrious}
\label{ch:bn-lazy-and-industrious}
\hlchaptermodality{462}

\begin{chapteropening}
\textbf{By the end of this chapter:} \emph{I can say whether someone is lazy, industrious or messy.}

Complete the last lesson of chapter 462: I can say whether someone is lazy, industrious or messy.
\end{chapteropening}

\section[\texorpdfstring{\bn{অলস} (ôlôs)}{অলস (ôlôs)}]{\bn{অলস} (ôlôs) — lazy; idle}

\label{lesson:BN-C462-olos}



\begin{quote}
Before the new one: say the Bengali for whimsical, then the Bengali for to dig one's heels in.
\end{quote}

\subsection*{You'll want to know: \bn{অলস}}
\textbf{\bn{অলস}} — \emph{ôlôs} — \textquotedblleft{}lazy; idle\textquotedblright{}.

\begin{quote}
\bn{ছুটির} \bn{দিনে} \bn{আমি} \bn{অলস} \bn{হয়ে} \bn{বিছানায়} \bn{শুয়ে} \bn{থাকি।}
\end{quote}

\begin{quote}
\emph{chhuṭir dine āmi ôlôs hôye bichhānāy shuye thāki.}
\end{quote}

\begin{quote}
On holidays I lie lazily in bed.
\end{quote}

\begin{cousinweb}
\textbf{\bn{অলস}} (\emph{ôlôs}) is a tatsama, Sanskrit \emph{alasa}, \textquotedblleft{}inactive, lazy\textquotedblright{} (CDIAL 708).
\end{cousinweb}

\begin{grammarlens}[title={What you've built}]
One more word for describing what you lived and felt.
\end{grammarlens}

\subsection*{Your turn}
\begin{itemize}
\raggedright
  \item \emph{Say it:} \emph{ôlôs}
  \item \emph{Say it:} \emph{ôlôs}, once more
  \item \emph{Say it:} \emph{chhuṭir dine āmi ôlôs hôye bichhānāy shuye thāki.} — On holidays I lie lazily in bed.
  \item \emph{Recall:} say the Bengali for whimsical, then the Bengali for to dig one's heels in, then say \emph{ôlôs} again
\end{itemize}

\begin{tcolorbox}[breakable,colback=teal!4,colframe=teal!35!black,title={Before you move on}]
What does \bn{অলস} mean? (\textquotedblleft{}Lazy; idle\textquotedblright{}.) Say it once more.
\end{tcolorbox}
\section[\texorpdfstring{\bn{কর্মঠ} (kôrmôṭh)}{কর্মঠ (kôrmôṭh)}]{\bn{কর্মঠ} (kôrmôṭh) — industrious; always busy at work}

\label{lesson:BN-C462-kormoth}



\begin{quote}
Before the new one: say the Bengali for to dig one's heels in, then the Bengali for lazy.
\end{quote}

\subsection*{You'll want to know: \bn{কর্মঠ}}
\textbf{\bn{কর্মঠ}} — \emph{kôrmôṭh} — \textquotedblleft{}industrious; always busy at work\textquotedblright{}.

\begin{quote}
\bn{দিদিমা} \bn{সারাদিন} \bn{কাজ} \bn{করেন}, \bn{তিনি} \bn{খুব} \bn{কর্মঠ।}
\end{quote}

\begin{quote}
\emph{didimā sārādin kāj kôren, tini khub kôrmôṭh.}
\end{quote}

\begin{quote}
My grandmother works all day long: she is very industrious.
\end{quote}

\begin{cousinweb}
\textbf{\bn{কর্মঠ}} (\emph{kôrmôṭh}) is a tatsama, Sanskrit \emph{karmaṭha}, \textquotedblleft{}capable of work, working diligently\textquotedblright{} (Monier-Williams).
\end{cousinweb}

\begin{grammarlens}[title={What you've built}]
One more word for describing what you lived and felt.
\end{grammarlens}

\subsection*{Your turn}
\begin{itemize}
\raggedright
  \item \emph{Say it:} \emph{kôrmôṭh}
  \item \emph{Say it:} \emph{kôrmôṭh}, once more
  \item \emph{Say it:} \emph{didimā sārādin kāj kôren, tini khub kôrmôṭh.} — My grandmother works all day long: she is very industrious.
  \item \emph{Recall:} say the Bengali for to dig one's heels in, then the Bengali for lazy, then say \emph{kôrmôṭh} again
\end{itemize}

\begin{tcolorbox}[breakable,colback=teal!4,colframe=teal!35!black,title={Before you move on}]
What does \bn{কর্মঠ} mean? (\textquotedblleft{}Industrious; always busy at work\textquotedblright{}.) Say it once more.
\end{tcolorbox}
\section[\texorpdfstring{\bn{এলোমেলো} (elomelo)}{এলোমেলো (elomelo)}]{\bn{এলোমেলো} (elomelo) — messy; in disorder}

\label{lesson:BN-C462-elomelo}



\begin{quote}
Before the new one: say the Bengali for lazy, then the Bengali for industrious.
\end{quote}

\subsection*{You'll want to know: \bn{এলোমেলো}}
\textbf{\bn{এলোমেলো}} — \emph{elomelo} — \textquotedblleft{}messy; in disorder\textquotedblright{}.

\begin{quote}
\bn{অরুণের} \bn{টেবিলটা} \bn{সবসময়} \bn{এলোমেলো} \bn{থাকে।}
\end{quote}

\begin{quote}
\emph{ôruner ṭebilṭā sôbsômôy elomelo thāke.}
\end{quote}

\begin{quote}
Arun's desk is always in a mess.
\end{quote}

\begin{cousinweb}
\textbf{\bn{এলোমেলো}} (\emph{elomelo}) is a rhyming doubling that paints things scattered every which way; its origin is not settled.
\end{cousinweb}

\begin{grammarlens}[title={What you've built}]
One more word for describing what you lived and felt.
\end{grammarlens}

\subsection*{Your turn}
\begin{itemize}
\raggedright
  \item \emph{Say it:} \emph{elomelo}
  \item \emph{Say it:} \emph{elomelo}, once more
  \item \emph{Say it:} \emph{ôruner ṭebilṭā sôbsômôy elomelo thāke.} — Arun's desk is always in a mess.
  \item \emph{Recall:} say the Bengali for lazy, then the Bengali for industrious, then say \emph{elomelo} again
\end{itemize}

\begin{tcolorbox}[breakable,colback=teal!4,colframe=teal!35!black,title={Before you move on}]
What does \bn{এলোমেলো} mean? (\textquotedblleft{}Messy; in disorder\textquotedblright{}.) Say it once more.
\end{tcolorbox}
\section[\texorpdfstring{\bn{ফাঁকি} \bn{দেওয়া} (phā̃ki deoyā)}{ফাঁকি দেওয়া (phā̃ki deoyā)}]{\bn{ফাঁকি} \bn{দেওয়া} (phā̃ki deoyā) — to shirk; to dodge (work)}

\label{lesson:BN-C462-phaki-deoya}



\begin{quote}
Before the new one: say the Bengali for industrious, then the Bengali for in disorder.
\end{quote}

\subsection*{You'll want to know: \bn{ফাঁকি} \bn{দেওয়া}}
\textbf{\bn{ফাঁকি} \bn{দেওয়া}} — \emph{phā̃ki deoyā} — \textquotedblleft{}to shirk; to dodge (work)\textquotedblright{}.

\begin{quote}
\bn{অলস} \bn{কর্মচারীটা} \bn{কাজে} \bn{ফাঁকি} \bn{দিত।}
\end{quote}

\begin{quote}
\emph{ôlôs kôrmôchārīṭā kāje phā̃ki ditô.}
\end{quote}

\begin{quote}
The lazy employee used to shirk his work.
\end{quote}

Past (he/she …): \bn{তিনি} \bn{কাজে} \bn{ফাঁকি} \bn{দিলেন}, \emph{tini kāje phā̃ki dilen}.

\begin{cousinweb}
\textbf{\bn{ফাঁকি} \bn{দেওয়া}} (\emph{phā̃ki deoyā}) is \emph{phā̃ki}, \textquotedblleft{}a trick, a dodge\textquotedblright{}, which CDIAL lists for Bengali (CDIAL 9035), with \emph{deoyā}, \textquotedblleft{}to give\textquotedblright{} (CDIAL 6141).
\end{cousinweb}

\begin{grammarlens}[title={What you've built}]
One more word for describing what you lived and felt.
\end{grammarlens}

\subsection*{Your turn}
\begin{itemize}
\raggedright
  \item \emph{Say it:} \emph{phā̃ki deoyā}
  \item \emph{Say it:} \emph{phā̃ki deoyā}, once more
  \item \emph{Say it:} \emph{ôlôs kôrmôchārīṭā kāje phā̃ki ditô.} — The lazy employee used to shirk his work.
  \item \emph{Recall:} say the Bengali for industrious, then the Bengali for in disorder, then say \emph{phā̃ki deoyā} again
\end{itemize}

\begin{tcolorbox}[breakable,colback=teal!4,colframe=teal!35!black,title={Before you move on}]
What does \bn{ফাঁকি} \bn{দেওয়া} mean? (\textquotedblleft{}To shirk; to dodge (work)\textquotedblright{}.) Say it once more.
\end{tcolorbox}
\section[\texorpdfstring{\bn{গড়িমসি} \bn{করা} (gôṛimôsi kôrā)}{গড়িমসি করা (gôṛimôsi kôrā)}]{\bn{গড়িমসি} \bn{করা} (gôṛimôsi kôrā) — to dawdle; to put off doing something}

\label{lesson:BN-C462-gorimosi-kora}



\begin{quote}
Before the new one: say the Bengali for in disorder, then the Bengali for to shirk.
\end{quote}

\subsection*{You'll want to know: \bn{গড়িমসি} \bn{করা}}
\textbf{\bn{গড়িমসি} \bn{করা}} — \emph{gôṛimôsi kôrā} — \textquotedblleft{}to dawdle; to put off doing something\textquotedblright{}.

\begin{quote}
\bn{বাড়ির} \bn{কাজ} \bn{করতে} \bn{ছেলেটা} \bn{গড়িমসি} \bn{করছিল।}
\end{quote}

\begin{quote}
\emph{bāṛir kāj kôrte chheleṭā gôṛimôsi kôrchhilô.}
\end{quote}

\begin{quote}
The boy was dawdling over his homework.
\end{quote}

Past (he/she …): \bn{তিনি} \bn{সই} \bn{করতে} \bn{গড়িমসি} \bn{করলেন}, \emph{tini sôi kôrte gôṛimôsi kôrlen}.

\begin{cousinweb}
\textbf{\bn{গড়িমসি} \bn{করা}} (\emph{gôṛimôsi kôrā}) is \emph{gôṛimôsi}, \textquotedblleft{}dawdling\textquotedblright{}, whose origin is not settled, with \emph{kôrā}, \textquotedblleft{}to do\textquotedblright{} (CDIAL 2814).
\end{cousinweb}

\begin{grammarlens}[title={What you've built}]
One more word for describing what you lived and felt.
\end{grammarlens}

\subsection*{Your turn}
\begin{itemize}
\raggedright
  \item \emph{Say it:} \emph{gôṛimôsi kôrā}
  \item \emph{Say it:} \emph{gôṛimôsi kôrā}, once more
  \item \emph{Say it:} \emph{bāṛir kāj kôrte chheleṭā gôṛimôsi kôrchhilô.} — The boy was dawdling over his homework.
  \item \emph{Recall:} say the Bengali for in disorder, then the Bengali for to shirk, then say \emph{gôṛimôsi kôrā} again
\end{itemize}

\begin{tcolorbox}[breakable,colback=teal!4,colframe=teal!35!black,title={Before you move on}]
What does \bn{গড়িমসি} \bn{করা} mean? (\textquotedblleft{}To dawdle; to put off doing something\textquotedblright{}.) Say it once more.
\end{tcolorbox}
